\documentclass[10pt,a4paper]{article}

\usepackage{fontawesome5}
\providecommand{\faPencilAlt}{\faIcon{pencil-alt}}

\usepackage{styles/couleurs}
\usepackage{styles/mise_en_page}
\usepackage{styles/boites}
\usepackage{styles/maths}
\usepackage{styles/figures}

\usepackage{enumitem}
\usepackage{amsmath}
\usepackage{amssymb}
\usepackage{tikz}
\usepackage{geometry}

\geometry{top=1cm,bottom=1cm,left=1.15cm,right=1.15cm}

\setlength{\parindent}{0pt}

\setlist[itemize]{
  leftmargin=1.35em,
  itemsep=0.18em,
  topsep=0.12em,
  parsep=0pt
}

\begin{document}

\begin{center}
  {\sffamily\LARGE\bfseries\color{TitreC}
  Le mode d'emploi pour rédiger en maths}\\[-0.05em]
  {\sffamily\small\color{SavoirC}
  À garder sous la main quand j'écris une solution}
\end{center}

\vspace{0.15em}

\begin{tcolorbox}[
  enhanced,
  colback=SavoirBg,
  colframe=SavoirC!65,
  boxrule=0.7pt,
  arc=2mm,
  left=3mm,right=3mm,top=1.5mm,bottom=1.5mm,
  before skip=1mm,after skip=2mm
]
\centering
{\sffamily\bfseries\color{SavoirC}
Mon objectif : qu'un camarade puisse comprendre ma solution dès la première lecture.}
\end{tcolorbox}

\vspace{0.2em}

% =========================================================
% DEUX COLONNES ALIGNÉES EN HAUT
% =========================================================

\begin{minipage}[t]{0.485\textwidth}
\vspace{0pt}

% =========================================================
% JE ME FAIS COMPRENDRE
% =========================================================

\begin{tcolorbox}[
  enhanced,
  colback=white,
  colframe=ActionC!65,
  boxrule=0.6pt,
  arc=2mm,
  title={\sffamily\bfseries Je me fais comprendre},
  coltitle=ActionC,
  fonttitle=\normalsize,
  left=2.5mm,right=2.5mm,top=1.4mm,bottom=1.4mm
]
\small
\begin{itemize}
  \item J'écris comme si j'expliquais à un camarade.

  \item Je ne raconte pas tout : j'écris ce qui aide vraiment à comprendre mon raisonnement.

  \item J'utilise les symboles et les notations mathématiques quand ils rendent l'écriture plus claire.

  \item Quand une phrase est nécessaire, je fais une vraie phrase.

  \item J'ai le droit de laisser dans ma rédaction un raisonnement qui n'aboutit pas,
  \textbf{sans l'effacer}, si j'indique clairement que j'abandonne cette piste.
\end{itemize}
\end{tcolorbox}

\vspace{0.3em}

% =========================================================
% CALCULS
% =========================================================

\begin{tcolorbox}[
  enhanced,
  colback=white,
  colframe=SavoirC!65,
  boxrule=0.6pt,
  arc=2mm,
  title={\sffamily\bfseries Mes calculs doivent rester vrais},
  coltitle=SavoirC,
  fonttitle=\normalsize,
  left=2.5mm,right=2.5mm,top=1.4mm,bottom=1.4mm
]
\small

\begin{itemize}
  \item Le signe \(=\) signifie \textbf{« a la même valeur que »}.
  Il ne veut pas dire « et puis ».

  \item Je peux écrire plusieurs égalités \textbf{sur la même ligne}
  si elles sont toutes vraies :
\end{itemize}

\vspace{-0.25em}

\[
24+6-5 = 30-5 = 25
\]

\vspace{-0.4em}

\begin{itemize}
  \item Je peux aussi écrire les étapes \textbf{les unes sous les autres} :
\end{itemize}

\vspace{-0.25em}

\[
\begin{aligned}
24+6-5 &= 30-5\\
       &= 25
\end{aligned}
\]

\vspace{-0.4em}

\begin{center}
\scriptsize\itshape
Les deux écritures sont correctes.
En colonne, il est souvent plus facile de suivre et de vérifier chaque étape.
\end{center}

\vspace{-0.25em}

\begin{itemize}
  \item Si un nombre n'était pas dans l'énoncé, je dois pouvoir expliquer d'où il vient par un calcul.

  \item Je peux séparer mes opérations si cela rend mon raisonnement plus clair.
\end{itemize}

\end{tcolorbox}

\end{minipage}
\hfill
\begin{minipage}[t]{0.485\textwidth}
\vspace{0pt}

% =========================================================
% JUSTIFIER AVEC UNE PROPRIÉTÉ
% =========================================================

\begin{tcolorbox}[
  enhanced,
  colback=AttentionC!5,
  colframe=AttentionC!70,
  boxrule=0.7pt,
  arc=2mm,
  title={\sffamily\bfseries Quand je justifie avec une propriété},
  coltitle=AttentionC,
  fonttitle=\normalsize,
  left=2.5mm,right=2.5mm,top=1.4mm,bottom=1.4mm
]

\centering
{\sffamily\bfseries\large
On sait que \(\ldots\)\\[0.25em]
Or \(\ldots\)\\[0.25em]
Donc \(\ldots\)}
\par

\vspace{0.4em}

\raggedright
\small
\begin{itemize}
  \item \textbf{On sait que} : j'écris les informations connues dans l'exercice.

  \item \textbf{Or} : j'écris la propriété, le théorème ou la règle que j'utilise.

  \item \textbf{Donc} : j'écris ce que j'en conclus.
\end{itemize}

\begin{center}
\scriptsize\itshape
Ce squelette m'aide à démarrer.
Je peux rédiger autrement si mon raisonnement reste clair.
\end{center}

\end{tcolorbox}

\vspace{0.3em}

% =========================================================
% BROUILLON
% =========================================================

\begin{tcolorbox}[
  enhanced,
  colback=white,
  colframe=NeutreC!55,
  boxrule=0.6pt,
  arc=2mm,
  title={\sffamily\bfseries Mon brouillon fait partie du travail},
  coltitle=TitreC,
  fonttitle=\normalsize,
  left=2.5mm,right=2.5mm,top=1.4mm,bottom=1.4mm
]
\small
\begin{itemize}
  \item Un brouillon n'est pas « sale » : il garde la trace de ma recherche.

  \item Une piste qui ne marche pas peut quand même m'aider.

  \item Le brouillon m'aide à \textbf{commencer ma rédaction} :
  je peux m'appuyer sur les étapes que j'ai déjà trouvées.

  \item J'évite de faire \textbf{recherche et rédaction en même temps},
  sauf si je suis sûre de moi.

  \item Je garde ce qui permet de comprendre comment j'ai trouvé.

  \item Avant de prendre la calculatrice, j’écris ce que je veux calculer ou tester.
\end{itemize}
\end{tcolorbox}

\end{minipage}

\vspace{0.35em}

% =========================================================
% CHECK-LIST FINALE
% =========================================================

\begin{tcolorbox}[
  enhanced,
  colback=NeutreC!4,
  colframe=NeutreC!40,
  boxrule=0.6pt,
  arc=2mm,
  left=3mm,right=3mm,top=1.4mm,bottom=1.4mm,
  before skip=1mm,after skip=0mm
]
{\sffamily\bfseries\color{TitreC}Avant de rendre, je vérifie :}

\vspace{0.25em}

\small
\(\square\) On comprend ce que j'ai fait.
\qquad
\(\square\) Mes signes \(=\) sont vrais.

\vspace{0.15em}

\(\square\) Les nombres nouveaux sont expliqués.
\qquad
\(\square\) J'ai justifié ce qui devait l'être.
\end{tcolorbox}

\end{document}